\chapter{Empirical Tests}
 
An empirical test for randomness is a procedure that takes a sequence of
observed draw outcomes and asks a single focused question about it.  Each
question probes a different property that a truly random sequence should
have --- such as whether all outcomes appear equally often, whether consecutive
draws are independent of each other, or whether patterns of repetition arise
at the right rate.
 
The word \emph{empirical} simply means that the test is based entirely on
the observed data.  No assumptions are made about how the draws were
produced.  The data speaks for itself.
 
\section{The Core Idea}
 
The central idea behind all empirical tests is the same:
 
\begin{enumerate}
  \item \textbf{Choose a property} that a random sequence should have.
  \item \textbf{Measure} how much the observed sequence agrees or disagrees
        with that property.
  \item \textbf{Decide} whether the disagreement is small enough to be
        attributed to natural fluctuation, or large enough to suggest a
        real problem.
\end{enumerate}
 
The measurement in step~2 always produces a single number --- the
\emph{test statistic}.  Step~3 then compares this number to the range of
values that would be typical for a genuinely random sequence.
 
\section{No Test Can Prove Randomness}
 
A crucial point that Knuth emphasises: no test, or even collection of tests,
can \emph{prove} that a sequence is random.  All a test can do is
\emph{fail to find} evidence of non-randomness in the specific property it
examines.
 
This means that passing a test is reassuring but not conclusive.  Failing a
test, on the other hand, is meaningful --- it tells you that something in the
sequence does not behave the way a fair process would.
 
\section{Many Tests Are Needed}
 
Because each test examines only one property, a sequence could look perfectly
fine from one angle while hiding a problem that only another angle would
reveal.  For example, a draw mechanism might produce each ball number the
correct number of times overall, yet still favour drawing the same number
twice in a row more often than chance allows.  The first property would pass
one test; the second would be caught by a different test.
 
For this reason, Knuth recommends applying a \emph{battery} of tests --- a
collection of tests each targeting a different property.  The more tests a
sequence passes, the more confidence we have in the fairness of the draw
mechanism.
